\documentclass[aspectratio=169,11pt]{beamer}

\usetheme{Madrid}
\usecolortheme{default}
\usefonttheme{professionalfonts}
\setbeamertemplate{navigation symbols}{}
\setbeamertemplate{footline}[frame number]

\usepackage{amsmath, amssymb, booktabs, array}

\title{Identity, Norms, Fairness, and Labor-Market Allocation}
\subtitle{Behavioral Labor -- Week 6}
\author{Teaching deck draft}
\date{}

\begin{document}

\begin{frame}
  \titlepage
\end{frame}

\begin{frame}{Central question}
\begin{itemize}
    \item How do identity, norms, fairness concerns, peer pressure, and firm culture shape who works where?
    \item Week 6 studies micro social objects in labor markets: roles, reference groups, teams, cultures, and supervisors.
    \item The labor margins are effort, morale, quits, promotion, pay satisfaction, applications, and worker-firm sorting.
    \item The discipline: name the social object, the margin, the reference group, and the identifying variation.
\end{itemize}
\end{frame}

\begin{frame}{Bridge from Weeks 1--5}
\small
\begin{tabular}{p{0.25\linewidth} p{0.60\linewidth}}
\toprule
Earlier object & Week 6 shift \\
\midrule
Behavioral taxonomy & Identity, norms, fairness, and culture become labor primitives \\
Nonstandard preferences & Social preferences and reference groups enter effort and job choice \\
Beliefs and attention & Workers form beliefs about what ``people like me'' do at work \\
Workplace contracts & Pay and monitoring are interpreted inside teams and firms \\
\bottomrule
\end{tabular}

\vspace{0.7em}
Week 5: contracts and supervision. Week 6: social meaning, micro norms, and allocation.
\end{frame}

\begin{frame}{Why this is behavioral labor}
\small
\begin{tabular}{p{0.25\linewidth} p{0.30\linewidth} p{0.29\linewidth}}
\toprule
Behavioral object & Labor margin & Empirical question \\
\midrule
Identity & job choice, effort, promotion & does role fit move behavior? \\
Fairness & morale, effort, quits & who is the reference group? \\
Peer pressure & productivity, attendance & learning or social sanction? \\
Firm culture & applications, retention & sorting or treatment? \\
Manager norms & promotion, pay gaps & quality or transmission? \\
\bottomrule
\end{tabular}
\end{frame}

\begin{frame}{Boundary with institutions course}
\small
\begin{tabular}{p{0.42\linewidth} p{0.42\linewidth}}
\toprule
Week 6: micro norms and allocation & Institutions course: macro/political rules \\
\midrule
teams, shifts, coworkers, supervisors & labor law and state capacity \\
fairness, comparison wages, morale & unions as political organizations \\
firm culture and job sorting & historical institutional change \\
manager transmission inside firms & macro cultural persistence \\
quits, effort, promotion, applications & political conflict over labor-market rules \\
\bottomrule
\end{tabular}

\vspace{0.7em}
If the norm works through workplace interaction or worker-firm sorting, it belongs here.
\end{frame}

\begin{frame}{Identity and self-image at work}
\[
U_i(a,j) = u\!\left(w_j\right) - c_i(a)
  + I_i\!\left(a, j, N_j, C_j\right)
\]

\small
\begin{tabular}{p{0.18\linewidth} p{0.68\linewidth}}
\toprule
Object & Interpretation \\
\midrule
$a$ & effort, compliance, application, promotion effort, or staying \\
$j$ & job, team, occupation, or firm \\
$N_j$ & local social norms in the job or team \\
$C_j$ & firm culture \\
$I_i(\cdot)$ & identity utility and role fit \\
\bottomrule
\end{tabular}
\end{frame}

\begin{frame}{Identity as a labor allocation object}
\begin{itemize}
    \item Identity can change the perceived cost or benefit of a job role.
    \item It can substitute for incentives when workers internalize role expectations.
    \item It can make incentives less effective when contracts conflict with self-image.
    \item It can shape occupational entry, firm sorting, promotion aspirations, and retention.
    \item Empirical discipline: identity object, labor margin, relevant group, and standard alternative.
\end{itemize}
\end{frame}

\begin{frame}{Fairness and pay comparisons}
\[
U_i = u\!\left(w_i\right) - c(e_i)
 - \phi_i \max\{0, w^{ref}_i - w_i\}
 - \psi_i \max\{0, w_i - w^{ref}_i\}
\]

\begin{itemize}
    \item $w^{ref}_i$ can be a peer wage, team wage, occupation wage, or manager wage.
    \item Fairness can move effort, attendance, morale, quits, cooperation, and pay satisfaction.
    \item Pay transparency may change fairness, beliefs about promotion, and outside search.
    \item Identification must separate relative-pay utility from information about career prospects.
\end{itemize}
\end{frame}

\begin{frame}{Peer pressure and coworker norms}
\[
U_i(e_i) = u\!\left(w_i\right) - c(e_i)
  - P_i\!\left(e_i,\bar e_{-i},N_g\right)
\]

\small
\begin{tabular}{p{0.22\linewidth} p{0.64\linewidth}}
\toprule
Object & Interpretation \\
\midrule
$\bar e_{-i}$ & coworker effort or visible peer behavior \\
$N_g$ & group norm \\
$P_i(\cdot)$ & social penalties from deviating from coworker behavior \\
Mas--Moretti & peers at work and productivity response \\
\bottomrule
\end{tabular}
\end{frame}

\begin{frame}{Learning or pressure?}
\small
\begin{tabular}{p{0.28\linewidth} p{0.27\linewidth} p{0.29\linewidth}}
\toprule
Channel & What changes? & Design clue \\
\midrule
Peer learning & information about productivity or technique & effects persist after exposure \\
Peer pressure & social cost of deviating & stronger when effort is visible \\
Productivity spillover & coworkers directly affect output & mechanical task complementarity \\
Common shock & shared demand or manager condition & same shift or team shock \\
Sorting & worker types choose peers or teams & exposure is endogenous \\
\bottomrule
\end{tabular}
\end{frame}

\begin{frame}{Firm culture and job sorting}
\[
j_i^\star \in \arg\max_{j \in \mathcal{J}}
 \; \mathbb{E}\!\left[U_i(j; X_i, N_j, C_j)\right]
\]

\begin{itemize}
    \item Culture affects applications, acceptance, retention, match quality, and mobility.
    \item Firms communicate culture through job postings, recruiting, managers, and rewards.
    \item Workers sort on expected norms and role fit.
    \item Huang--Pacelli--Shi--Zou: culture communication in labor markets.
\end{itemize}
\end{frame}

\begin{frame}{Sorting, treatment, and selection}
\small
\begin{tabular}{p{0.25\linewidth} p{0.32\linewidth} p{0.27\linewidth}}
\toprule
Claim & What must move? & Main caution \\
\midrule
Sorting & applicant or hire composition & culture text may reveal firm type \\
Treatment & incumbent behavior after exposure & workers may already fit culture \\
Retention & quits and tenure & exits change composition \\
Transmission & norms move through people & manager quality confounding \\
Equilibrium selection & worker and firm distributions shift & partial-equilibrium evidence \\
\bottomrule
\end{tabular}
\end{frame}

\begin{frame}{Managers and micro culture}
\[
C_{f,t+1} = \Gamma\!\left(C_{f,t}, M_{f,t}, H_{f,t}, R_{f,t}\right)
\]

\begin{itemize}
    \item Managers are not only monitors or evaluators.
    \item They transmit norms about effort, flexibility, ambition, belonging, and promotion.
    \item Culture evolves with managers, hiring, and reward systems.
    \item Frontier design: manager rotation or assignment plus worker outcomes.
\end{itemize}
\end{frame}

\begin{frame}{Empirical designs}
\scriptsize
\begin{tabular}{p{0.29\linewidth} p{0.28\linewidth} p{0.28\linewidth}}
\toprule
Design & Object identified & Interpretation problem \\
\midrule
Peer exposure field experiment & coworker norms and pressure & pressure versus learning \\
Pay-comparison experiment & fairness and reference wages & fairness versus career beliefs \\
Survey plus admin data & wage beliefs, morale, quits & stated beliefs versus behavior \\
Job-posting text design & culture communication & sorting versus firm type \\
Manager assignment design & norm transmission & manager quality versus culture \\
\bottomrule
\end{tabular}
\end{frame}

\begin{frame}{Week 6 lab}
\begin{itemize}
    \item Reproduce: synthetic peer-at-work factbook anchored on Mas--Moretti.
    \item Diagnose: norm or identity object, reference group, labor margin, identifying variation.
    \item Transfer: fairness and pay comparison anchored on Breza--Kaur--Shamdasani.
    \item Frontier extension: culture communication or manager-driven norm transmission.
\end{itemize}
\end{frame}

\begin{frame}{Frontier directions}
\begin{itemize}
    \item Measuring culture with postings, reviews, communications, and manager histories.
    \item Pay transparency with rich morale, effort, quit, and promotion outcomes.
    \item Algorithmic teams and digital peer pressure.
    \item Manager rotation as cultural transmission, not only supervision.
    \item Linking micro culture to worker-firm reallocation and inequality.
\end{itemize}
\end{frame}

\begin{frame}{Bridge to Week 7}
\begin{itemize}
    \item Week 6 names the social objects: identity, fairness, peer pressure, culture, manager norms.
    \item The next challenge is identification.
    \item Week 7 asks how to distinguish behavioral wedges from constraints, institutions, sorting, and measurement error.
    \item The recurring rule: object, margin, reference group, variation, interpretation.
\end{itemize}
\end{frame}

\begin{frame}{Takeaways}
\begin{itemize}
    \item Identity and norms matter when they move labor allocation and workplace behavior.
    \item Fairness and pay comparisons require explicit reference groups.
    \item Peer pressure, peer learning, and productivity spillovers are different mechanisms.
    \item Firm culture affects sorting, retention, and match quality.
    \item Managers can transmit norms, but credible evidence must separate transmission from assignment and quality.
\end{itemize}
\end{frame}

\end{document}
